\section{matching between lattice and continuum through lattice perturbation theory}
In this section, we aim at establishing the one-loop matching connecting the quasi quark distribution on the lattice to the normal quark distribution in the continuum. On the lattice, besides the diagrams in Fig.~\ref{1loopFeyn}, there exists an extra one-loop diagram (shown in Fig.~\ref{1loopLQCD}). 

The first diagram in Fig.~\ref{1loopFeyn} and the diagram in Fig.~3 are nothing but the Feynman gauge quark field renormalization, and have been well understood in lattice QCD. Their contribution can also be obtained from the other diagrams in Fig.~\ref{1loopFeyn} through quark number conservation. As shown in Refs.~\cite{Xiong:2013bka} and~\cite{Ji:2015jwa}, in the continuum the second and third diagrams in Fig.~\ref{1loopFeyn} contain a logarithmic dependence on the hadron momentum instead of a logarithmic divergence. Their contribution is independent of the regularization and thus will be the same in the continuum and on the lattice. One only needs to  match the linearly divergent part of the last diagram in Fig.~\ref{1loopFeyn} between lattice and continuum. This gives a connection between the lattice cutoff and the linear divergence in the one-loop kernel in the continuum~\cite{Xiong:2013bka} (in general they will differ by a finite term). %After matching the continuum linear divergence and the lattice cutoff, we can able to remove  %The lattice cutoff will be renormalized by the mass counterterm introduced in the previous section. Taking into account the remaining finite term, the one-loop kernel in the continuum can then be used to connect the lattice quasi distribution and the normal distribution. To realize this, we need to compute the Wilson line self energy using lattice perturbation theory.
To obtain this connection, we compute the Wilson line self energy using lattice perturbation theory based on a simple version of discretized gauge action.


\begin{figure}[tbp]
\centering
\includegraphics[width=0.2\textwidth]{images/1loopL1}
\caption{The extra one-loop diagram for quasi quark distribution on the lattice, the conjugate diagram is not shown.}
\label{1loopLQCD}
%\vspace*{-.5em}
\end{figure}


The contribution of the Wilson line self energy diagram can be written as
\begin{align}
\delta W&=\frac{\alpha_s C_F}{8\pi^3}\int d^4(k a_L) \frac{1-\cos(k_z z)}{(1-\cos(k_z a_L))4\sum_\lambda \sin^2\frac{k_\lambda a_L}{2}}\non\\
&=\alpha_s C_F[\frac{\pi}{2}\frac{|z|}{a_L}-\frac{1}{\pi}\ln\frac{\pi |z|}{a_L}-\frac{1}{\pi}]+\mathcal O(\frac{a_L}{z}),
\end{align}
where $a_L$ denotes the lattice spacing. Comparing this with the expansion of Eq.~(\ref{WLcoordspace}) leads to the following matching between the linear divergence in the continuum and on the lattice
\beq
\Lambda=\frac{\pi}{a_L}.
\eeq
By replacing $\Lambda$ in the one-loop kernel of Ref.~\cite{Xiong:2013bka} with the above matching condition, we then obtain the one-loop kernel relating the quasi PDF on the lattice to the normal PDF. By adding a factor of $e^{-\delta m |z|}$ to the integrand of Eq.~(\ref{qUnpolDef}), we obtain the improved quasi-PDF $\tilde q_{\text{imp}}$ free of power divergence
\beq\label{im}
   \tilde q_{\text{imp}}(x, \Lambda, p^z) = \int^\infty_{-\infty} \frac{dz}{4\pi} e^{izk^z-\delta m |z|}  \langle p|\overline{\psi}(0, 0_\perp, z)
   \gamma^z L(z,0)\psi(0) |p\rangle \ .
\eeq

The matching kernel $Z$ between the improved quasi-PDF $\tilde q_{\text{imp}}$ and the PDF $q$ is defined as
\beq
\tilde q_{\text{imp}}(x, a_L, p^z)  = \int^1_{-1} \frac{dy}{|y|} Z\left(\frac{x}{y},{p^z a_L},\frac{\mu}{p^z}\right) q(y, \mu) +  {\cal O}\left(\Lambda^2_{\rm QCD}/(p^z)^2,  M^2/(p^z)^2\right)\ .
\eeq
At one-loop the $Z$ factor can be extracted from Eqs.~(15)-(19) of Ref.~\cite{Xiong:2013bka} with the $1/a_L$ contribution 
subtracted by the counterterm diagram in Fig.~\ref{1loopFeynmassCT}.
\begin{equation}
             Z\left(\xi,  p^z a_L,\frac{\mu}{p^z}\right) = \delta (\xi-1) + \frac{\alpha_s}{2\pi} Z^{(1)}
             \left(\xi, p^z a_L,\frac{\mu}{p^z}\right) + \dots ,
\end{equation}
where
\begin{align}
Z^{(1)}/C_F
=\left\{ \begin{array} {ll} \left(\frac{1+\xi^2}{1-\xi}\right)\ln \frac{\xi}{\xi-1} + 1\ , & \xi>1\ , \\ \left(\frac{1+\xi^2}{1-\xi}\right)\ln \frac{(p^z)^2}{\mu^2}+\left(\frac{1+\xi^2}{1-\xi}\right)\ln\big[{4\xi(1-\xi)}\big]-\frac{2\xi}{1-\xi}+1\ , & 0<\xi<1\ , \\ \left(\frac{1+\xi^2}{1-\xi}\right)\ln \frac{\xi-1}{\xi}-1\ , & \xi<0\ , \end{array} \right.
\end{align}
and the extra contribution near $\xi=1$ is
\begin{align}
\delta Z^{(1)}
=\frac{\alpha_S C_F}{2\pi} \int dy \left\{ \begin{array} {ll} -\frac{1+y^2}{1-y}\ln \frac{y}{y-1}-1\ , & y>1\ , \\ -\frac{1+y^2}{1-y}\ln \frac{(p^z)^2}{\mu^2}-\frac{1+y^2}{1-y}\ln\big[4y(1-y)\big]+\frac{2y(2y-1)}{1-y}+1\ , & 0<y<1\ , \\ -\frac{1+y^2}{1-y}\ln \frac{y-1}{y}+1\ , & y<0\ . \end{array} \right.
\end{align}
Given that the power divergence has been removed, the above one-loop matching kernel facilitates a reliable extraction of normal PDFs from lattice data. 


